\documentclass[11pt,a4paper]{article}
\usepackage{fontspec}
\usepackage[margin=18mm]{geometry}
\usepackage{graphicx}
\usepackage{booktabs}
\usepackage{tabularx}
\usepackage{array}
\usepackage{float}
\usepackage{caption}
\usepackage{hyperref}
\usepackage{xcolor}

\defaultfontfeatures{Ligatures=TeX,Scale=MatchLowercase}
\IfFontExistsTF{Noto Sans CJK KR}{\setmainfont{Noto Sans CJK KR}\setsansfont{Noto Sans CJK KR}}{\setmainfont{Apple SD Gothic Neo}\setsansfont{Apple SD Gothic Neo}}
\IfFontExistsTF{DejaVu Sans Mono}{\setmonofont{DejaVu Sans Mono}}{\setmonofont{Menlo}}
\XeTeXlinebreaklocale "ko"
\XeTeXlinebreakskip = 0pt plus 1pt
\hypersetup{colorlinks=true,linkcolor=blue!50!black,urlcolor=blue!50!black}
\setlength{\parskip}{0.45em}
\setlength{\parindent}{0pt}
\emergencystretch=2em
\renewcommand{\arraystretch}{1.16}
\newcolumntype{Y}{>{\raggedright\arraybackslash}X}
\newcommand{\code}[1]{\texttt{#1}}
\captionsetup{font=small,labelfont=bf,justification=raggedright,singlelinecheck=false}

\title{2026-04-01 UUV MuJoCo A/B 원인분리 분석\\\large inertia-only vs thruster-gain-only}
\author{}
\date{2026-04-28}

\begin{document}
\maketitle

\section{A/B 구성}
\begin{figure}[H]
  \centering
  \includegraphics[width=0.98\linewidth]{real_bag_2026_04_01_ab_comparison/latex_figures/ab_pipeline.png}
  \caption{baseline, inertia-only, thruster-gain-only, patched-current 비교 구조. 비교 파이프라인과 sensor frame correction은 고정했다. patched-current는 실제 \code{uuv\_mujoco} config 수정 후 재실행한 결과다.}
\end{figure}

\section{결과 그래프}
\begin{figure}[H]
  \centering
  \includegraphics[width=0.98\linewidth]{real_bag_2026_04_01_ab_comparison/latex_figures/ab_rmse.png}
  \caption{RMSE 비교. thruster-gain-only는 DVL/gyro/depth RMSE를 가장 크게 줄인다. inertia-only는 일부 RMSE를 낮추지만 상관/gain 개선이 일관되지 않다.}
\end{figure}

\begin{figure}[H]
  \centering
  \includegraphics[width=0.98\linewidth]{real_bag_2026_04_01_ab_comparison/latex_figures/ab_corr_gain.png}
  \caption{상관계수와 real/sim gain. \code{gain\_scale\_all=1.0}은 yaw/DVL amplitude gap을 줄이는 방향으로 작동하지만, 아직 gain 1.0까지는 부족하다. patched-current는 combined override와 거의 같은 결과다.}
\end{figure}

\begin{figure}[H]
  \centering
  \includegraphics[width=0.94\linewidth]{real_bag_2026_04_01_ab_comparison/latex_figures/ab_thruster_contact.png}
  \caption{yaw thruster force p95와 contact fraction. thruster gain 수정 후 yaw p95 force는 약 149 N에서 51 N 수준으로 낮아진다. contact fraction은 여전히 높아서 depth 평가는 controller-in-loop로 다시 해야 한다.}
\end{figure}

\section{판단}
\begin{table}[H]
  \centering
  \caption{핵심 gain 비교}
  \begin{tabularx}{\linewidth}{lYYYY}
    \toprule
    Bag/Case & DVL x gain & DVL y gain & gyro z gain & 해석 \\
    \midrule
    20:08 baseline & 0.254 & 0.150 & 0.351 & sim amplitude 과대 \\
    20:08 inertia-only & 0.245 & 0.144 & 0.358 & 관성만으로 개선 안 됨 \\
    20:08 thruster-gain-only & 0.377 & 0.211 & 0.573 & 가장 큰 개선 \\
    20:08 patched-current & 0.376 & 0.212 & 0.608 & 실제 config 반영 후 결과 \\
    20:20 baseline & 0.082 & 0.103 & 0.157 & sim amplitude 매우 과대 \\
    20:20 inertia-only & 0.027 & 0.097 & 0.134 & 오히려 일부 악화 \\
    20:20 thruster-gain-only & 0.119 & 0.145 & 0.194 & 개선되나 아직 부족 \\
    20:20 patched-current & 0.130 & 0.197 & 0.203 & 실제 config 반영 후 결과 \\
    \bottomrule
  \end{tabularx}
\end{table}

결론: 현재 핵심 문제는 \code{body\_inertia\_scale\_xyz}보다 \code{gain\_scale\_all=2.9} 쪽 신호가 더 강하다. 실제 patch는 \code{gain\_scale\_all=1.0}과 \code{body\_inertia\_scale\_xyz=[1,1,1]}를 반영했다. 관성 정상화는 물리적으로 필요하지만 단독으로 bag fit을 해결하지 못했고, 추력 gain 수정이 RMSE/gain 개선의 주요 원인이다.

\end{document}